% ============================================================
% 理论力学 · 交互讲义 —— M7《刚体几何与角动量》note.tex
% ============================================================
\documentclass[UTF8,a4paper,11pt]{ctexart}
\usepackage{amsmath,amssymb}
\usepackage{geometry}
\usepackage{booktabs}
\usepackage{hyperref}
\geometry{margin=2.5cm}
\hypersetup{colorlinks=true,linkcolor=blue,urlcolor=blue}

\title{\textbf{刚体几何与角动量}}
\author{理论力学 · 交互讲义}
\date{\today}

\begin{document}
\maketitle

\begin{abstract}
本讲义处理刚体几何与角动量：球冠质心、甲虫-圆环角动量守恒、
圆盘角动量、Newton 球壳定理。全部解析结果由配套脚本
\texttt{qa/numeric\_check.py} 逐项数值验证。本模块覆盖作业-4。
\end{abstract}

\tableofcontents
\newpage

% ============================================================
\section{作业映射}
\label{sec:map}

\begin{center}
\small
\begin{tabular}{cp{6.5cm}p{3.5cm}c}
\toprule
来源 & 题面简述 & 核心知识点 & 对应小节 \\
\midrule
作业-4 Q1 & 球冠质心（球冠高度 $h$） &
  质心积分、几何 & \ref{sub:spherical-cap} \\
作业-4 Q2 & 甲虫在光滑圆环上爬行，环转角 &
  角动量守恒、约束 & \ref{sub:beetle} \\
作业-4 Q3 & 薄圆盘绕中心轴角动量 &
  转动惯量、角动量 & \ref{sub:disk} \\
作业-4 Q4 & 球壳对质点的引力势能 &
  Newton 球壳定理 & \ref{sub:shell} \\
\bottomrule
\end{tabular}
\end{center}

% ============================================================
\section{模型与假设}
\label{sec:model}

\begin{enumerate}
  \item \textbf{刚体}：物体为刚体，形状不变。
  \item \textbf{质点/连续体}：按问题需要取质点或连续体模型。
  \item \textbf{光滑约束}：圆环光滑，无摩擦。
  \item \textbf{单位制}：SI 单位。
\end{enumerate}

% ============================================================
\section{推导}
\label{sec:derive}

\subsection{球冠质心（作业-4 Q1）}
\label{sub:spherical-cap}

球半径 $R$，切割平面距球心 $b$（$0\le b\le R$），球冠高度 $h=R-b$。
球冠体积
\begin{equation}
  V=\int_{b}^{R}\pi(R^2-z^2)\,\mathrm dz
  =\pi\left[R^2(R-b)-\frac{R^3-b^3}{3}\right]
  =\frac{\pi(R-b)^2(2R+b)}{3}.
\end{equation}
质心 $z$ 坐标（从球心）：
\begin{equation}
  z_{\rm cm}=\frac{1}{V}\int_b^R z\,\pi(R^2-z^2)\,\mathrm dz,
  \qquad
  \int_b^R z(R^2-z^2)\,\mathrm dz
  =\frac{(R^2-b^2)^2}{4},
\end{equation}
代入体积得
\begin{equation}
  z_{\rm cm}=\frac{3(R+b)^2}{4(2R+b)}.
\end{equation}
质心距底面（平面 $z=b$）的距离：
\begin{equation}
  d=z_{\rm cm}-b
  =\frac{3(R+b)^2}{4(2R+b)}-b
  =\frac{(R-b)(3R+b)}{4(2R+b)}.
\end{equation}
用球冠高度 $h=R-b$ 表示，即
\begin{equation}
  \boxed{\;d=\frac{h(4R-h)}{4(3R-h)}\;}
\end{equation}
当 $h=R$（半球）：$d=3R/8$；当 $h=0$（无球冠）：$d=0$。

\subsection{甲虫与圆环（作业-4 Q2）}
\label{sub:beetle}

圆环质量 $M$，半径 $R$，甲虫质量 $m$。圆环放在光滑水平面上，系统初始静止。
设甲虫相对圆环的方位角为 $\theta$，圆环自转角为 $\phi$，甲虫相对环心的实验室方位角为
\begin{equation}
  \psi=\phi+\theta.
\end{equation}
水平方向无外力，系统质心固定。取质心为原点，环心与甲虫到质心的距离分别为
\begin{equation}
  a_C=\frac{mR}{M+m},\qquad a_B=\frac{MR}{M+m},
\end{equation}
且
\begin{equation}
  \mathbf r_C=-a_C(\cos\psi,\sin\psi),\qquad
  \mathbf r_B=+a_B(\cos\psi,\sin\psi).
\end{equation}
关于质心的总角动量守恒。薄环绕自身中心的转动惯量为 $MR^2$；环心平移和甲虫绕质心转动分别为
$M a_C^2\dot\psi$、$m a_B^2\dot\psi$：
\begin{equation}
  L=MR^2\dot\phi+\left(Ma_C^2+ma_B^2\right)\dot\psi=0.
\end{equation}
其中
\begin{equation}
  Ma_C^2+ma_B^2
  =M\frac{m^2R^2}{(M+m)^2}+m\frac{M^2R^2}{(M+m)^2}
  =\frac{MmR^2}{M+m}.
\end{equation}
因此
\begin{equation}
  \dot\phi=-\frac{m}{M+m}\dot\psi.
\end{equation}
代入 $\psi=\phi+\theta$：
\begin{equation}
  \dot\phi=-\frac{m}{M+2m}\dot\theta,\qquad
  \dot\psi=\frac{M+m}{M+2m}\dot\theta.
\end{equation}
甲虫相对环爬完一圈 $\theta=2\pi$ 时，
\begin{equation}
  \boxed{\;\Delta\phi=-\frac{2\pi m}{M+2m},\qquad
  \Delta\psi=\frac{2\pi(M+m)}{M+2m}\;}
\end{equation}
即环反转 $2\pi m/(M+2m)$，甲虫在实验室系中转过 $2\pi(M+m)/(M+2m)$；质心始终固定，环心绕质心平移。
当 $M=m$ 时恢复 $120^\circ/240^\circ$。若把环心固定（外部约束），才会得到环转 $180^\circ$，但那时质心并不固定。

\paragraph{与均匀实心球的对比。}
均匀实心球内部
\begin{equation}
  V_{\rm solid}(r)=-\frac{GMm}{2a}\left(3-\frac{r^2}{a^2}\right)\quad(r<a),
\end{equation}
外部仍为 $-GMm/r$。球壳内部 $V$ 为常数；实心球内部随 $r$ 二次下降，
中心处比球壳表面低 $GMm/(2a)$。两者在 $r=a$ 连续，且 $r>a$ 完全相同。

\subsection{薄圆盘角动量（作业-4 Q3）}
\label{sub:disk}

薄圆盘质量 $M$，半径 $R$，绕中心垂直轴以角速度 $\omega_0$ 转动。
转动惯量 $I=\frac12 MR^2$。角动量
\begin{equation}
  \boxed{\;L=I\omega_0=\frac12 MR^2\omega_0\;}
\end{equation}

\subsection{Newton 球壳定理（作业-4 Q4）}
\label{sub:shell}

质量为 $M$ 的均匀球壳，半径 $a$。质点质量 $m$ 在距中心 $r$ 处。取无穷远处势能为 0，即 $V(\infty)=0$。

\paragraph{壳内（$r<a$）：} 引力势能为常数：
\begin{equation}
  V=-\frac{GMm}{a}.
\end{equation}
引力为零。

\paragraph{壳外（$r>a$）：} 等效于质量 $M$ 集中于中心：
\begin{equation}
  V=-\frac{GMm}{r}.
\end{equation}

\paragraph{球壳表面（$r=a$）：} $V=-GMm/a$，与壳内连续。

% ============================================================
\section{解析结果}
\label{sec:result}

\begin{enumerate}
  \item 球冠质心：$d=h(4R-h)/[4(3R-h)]$（\eqref{sub:spherical-cap}，$h$ 为球冠高度）。
  \item 甲虫-圆环（自由环）：环反转 $2\pi m/(M+2m)$，甲虫实验室方位转过 $2\pi(M+m)/(M+2m)$；$M=m$ 时为 $120^\circ/240^\circ$，质心固定（\eqref{sub:beetle}）。
  \item 圆盘角动量：$L=\frac12 MR^2\omega_0$（\eqref{sub:disk}）。
  \item Newton 球壳：$r<a$ 时 $V=-GMm/a$（常数），$r>a$ 时 $V=-GMm/r$
        （\eqref{sub:shell}）。
\end{enumerate}

% ============================================================
\section{数值验证}
\label{sec:numeric}

全部数值实验由配套脚本 \texttt{qa/numeric\_check.py} 完成（仅依赖 numpy）。

\begin{table}[htbp]
\centering
\small
\begin{tabular}{p{5cm}cc}
\toprule
物理量 & 解析值 & 数值值 \\
\midrule
球冠质心（$R=1,h=R$，半球） & $3/8=0.375$ & $0.375$ \\
球冠质心（$R=1,h=0.5$） & $0.5\cdot3.5/(4\cdot2.5)=0.175$ & $0.175$ \\
甲虫环反转角（M=m=1） & $2\pi/3=2.09440$ & $2.09440$ \\
圆盘 $L$（$M=1,R=1,\omega=1$） & $0.5$ & $0.5$ \\
球壳 $V$（$r=0.5a$） & $-GMm/a$ & $-GMm/a$ \\
球壳 $V$（$r=2a$） & $-GMm/(2a)$ & $-GMm/(2a)$ \\
\bottomrule
\end{tabular}
\caption{数值验证汇总。全部数值取自 \texttt{qa/numeric\_check.py} 的实际输出。}
\label{tab:numeric}
\end{table}

% ============================================================
\section{转动惯量与平行轴定理}
\label{sec:parallel-axis}

平行轴定理：$I=I_{\rm cm}+Md^2$，其中 $d$ 为新轴到质心轴的距离。
对圆盘：$I_{\rm cm}=\frac12 MR^2$（中心轴），$I_{\rm edge}=I_{\rm cm}+MR^2=\frac32 MR^2$（边缘轴）。

\section{角动量守恒的应用}
\label{sec:angular-momentum-apps}

\paragraph{滑冰者加速。}
滑冰者收臂时转动惯量减小，角速度增大（角动量守恒）。

\paragraph{地球-月球系统。}
潮汐摩擦导致地球自转减慢，月球轨道半径增大（角动量守恒）。

\section{Newton 球壳定理的证明概要}
\label{sec:shell-proof}

球壳定理的证明利用高斯定律类比：引力场 $\boldsymbol g$ 的通量
$\oint \boldsymbol g\cdot\mathrm d\boldsymbol A=-4\pi GM_{\rm enc}$。
壳内 $M_{\rm enc}=0$，故 $\boldsymbol g=0$；壳外 $M_{\rm enc}=M$，
故 $\boldsymbol g=-GM/r^2\,\boldsymbol e_r$。

\section{转动惯量的计算方法}
\label{sec:moment-inertia}

转动惯量 $I=\int r^2\,\mathrm dm$。常见形状：
\begin{itemize}
  \item 质点：$I=mr^2$。
  \item 薄环：$I=MR^2$。
  \item 薄盘：$I=\frac12 MR^2$。
  \item 实心球：$I=\frac25 MR^2$。
  \item 球壳：$I=\frac23 MR^2$。
\end{itemize}

\section{刚体动力学的历史}
\label{sec:history-rigid}

从 Euler 的刚体运动方程（1750s）到 Lagrange 的分析力学，
刚体动力学的发展经历了从几何到分析的转变。
Euler 方程 $\boldsymbol I\dot{\boldsymbol\omega}+\boldsymbol\omega\times\boldsymbol I\boldsymbol\omega=\boldsymbol M$
描述了刚体在惯性系中的转动。

\paragraph{Euler 角。}
三个 Euler 角（$\phi,\theta,\psi$）描述刚体取向，
是旋转矩阵的参数化。角速度在体坐标系中的分量为
$\omega_1=\dot\theta\cos\psi+\dot\phi\sin\theta\sin\psi$ 等。

\paragraph{陀螺运动。}
重陀螺的进动与章动是 Euler 方程的经典解，
在导航（陀螺仪）中有重要应用。

\section{附录：常见形状的转动惯量}
\label{sec:appendix-inertia}

\begin{table}[htbp]
\centering
\small
\begin{tabular}{p{4cm}c}
\toprule
形状 & $I$（中心轴） \\
\midrule
质点 & $mr^2$ \\
细杆（中心） & $\frac{1}{12}ml^2$ \\
细杆（端点） & $\frac{1}{3}ml^2$ \\
薄环 & $MR^2$ \\
薄盘 & $\frac12 MR^2$ \\
实心球 & $\frac25 MR^2$ \\
球壳 & $\frac23 MR^2$ \\
圆柱（轴） & $\frac12 MR^2$ \\
\bottomrule
\end{tabular}
\caption{常见形状的转动惯量。}
\label{tab:inertia}
\end{table}

\section{讨论}
\label{sec:discuss}

\paragraph{角动量守恒的普遍性。}
甲虫-圆环问题是角动量守恒的经典演示。系统初始角动量为零，
甲虫运动时环反向旋转，保持总角动量为零。

\paragraph{Newton 球壳定理的意义。}
球壳定理是 Newton 万有引力理论的基石，解释了为什么行星可视为质点，
以及为什么球壳内部引力为零。

\paragraph{球冠质心的应用。}
球冠质心公式用于计算球形天体（如行星、恒星）的引力势与转动惯量。

\paragraph{Newton 球壳定理的应用。}
球壳定理解释了为什么行星内部（r<R）引力随 r 线性增加（$g\propto r$），
而外部（r>R）按 $1/r^2$ 衰减。

\paragraph{与作业题的对照。}
表~\ref{tab:numeric} 覆盖了作业-4 Q1--Q4 的全部数值结果。

% ============================================================
\begin{thebibliography}{9}

\bibitem{goldstein}
H.~Goldstein, C.~P.~Poole, J.~L.~Safko,
\emph{Classical Mechanics}, 3rd ed., Pearson, 2002.
（刚体、角动量、引力势。）

\bibitem{newton}
I.~Newton, \emph{Principia Mathematica}, 1687.
（球壳定理。）

\bibitem{zhou}
周衍柏，《理论力学教程》（第三版），高等教育出版社，2009.
（刚体几何与角动量。）

\end{thebibliography}

\end{document}
